\section{2012 Geometry \& Topology}

\clearpage

\subsection{Individual Exam}

\begin{exercise}
\label{geometry:2012:1}
Show that $\pi_{3}(S^{2}) \neq 0$.
\end{exercise}

\begin{solution}
\textbf{Geometric Intuition:} In low dimensions, we often expect $\pi_n(S^k) = 0$ for $n > k$ because a lower-dimensional sphere should "rattle around" inside a higher-dimensional one and be shrinkable. However, the Hopf fibration\sidenote{The Hopf fibration can be constructed by identifying $S^3$ with the unit sphere in $\mathbb{C}^2$, i.e., $\{(z_1, z_2) \in \mathbb{C}^2 : |z_1|^2 + |z_2|^2 = 1\}$, and $S^2$ with the complex projective line $\mathbb{CP}^1$. The map $h: S^3 \to S^2$ is then the projection $h(z_1, z_2) = [z_1:z_2]$, which sends a point on the 3-sphere to the complex line through it and the origin.} $S^3 \to S^2$ shows that $S^3$ can be non-trivially "wrapped" around $S^2$ such that the inverse image of every point is a linked circle ($S^1$). This "linking" is what prevents the map from being null-homotopic.

\textbf{Proof:} Consider the Hopf fibration $S^1 \hookrightarrow S^3 \xrightarrow{h} S^2$. This is a fiber bundle with base $S^2$, total space $S^3$, and fiber $S^1$. The long exact sequence of homotopy groups for a fibration gives:
\[ \dots \to \pi_3(S^1) \to \pi_3(S^3) \to \pi_3(S^2) \to \pi_2(S^1) \to \dots \]
It is well-known that $\pi_n(S^1) = 0$ for $n > 1$. Thus $\pi_3(S^1) = 0$ and $\pi_2(S^1) = 0$. The sequence becomes:
\[ 0 \to \pi_3(S^3) \xrightarrow{h_*} \pi_3(S^2) \to 0 \]
This implies that $h_*: \pi_3(S^3) \to \pi_3(S^2)$ is an isomorphism. Since $\pi_3(S^3) \cong \mathbb{Z}$ (generated by the identity map), we conclude that $\pi_3(S^2) \cong \mathbb{Z} \neq 0$. The generator is the homotopy class of the Hopf map $h$.
\end{solution}

\begin{remark}[The Homotopy Groups of Spheres: A Deeper Look]
The computation of $\pi_n(S^k)$ is a central, notoriously difficult problem in algebraic topology. No general formula exists, and the known results reveal a rich, complex structure.

\textbf{1. The Simple Cases (Recap):}
\begin{itemize}
    \item For $n < k$, $\pi_n(S^k) = 0$.
    \item For $n = k$, $\pi_k(S^k) \cong \mathbb{Z}$ for $k \ge 1$ (via degree).
    \item For $k = 1$, $\pi_1(S^1) \cong \mathbb{Z}$ and $\pi_n(S^1)=0$ for $n>1$.
\end{itemize}

\textbf{2. From Unstable to Stable: The Hopf Map Example}

You correctly noted that the relationship between $\pi_3(S^2) \cong \mathbb{Z}$ and the stable group $\pi_1^s \cong \mathbb{Z}_2$ is confusing. This is a key point that illustrates the transition from the "unstable" to the "stable" range.
\begin{itemize}
    \item The map $h: S^3 \to S^2$ (Hopf map) generates the unstable group $\pi_3(S^2) \cong \mathbb{Z}$.
    \item The \textbf{suspension homomorphism} $E: \pi_n(S^k) \to \pi_{n+1}(S^{k+1})$ is defined by taking a map $f: S^n \to S^k$ and extending it to a map on the suspensions $S(f): S(S^n) \to S(S^k)$. Since $S(S^m) \cong S^{m+1}$, this gives a map from $\pi_n(S^k)$ to $\pi_{n+1}(S^{k+1})$.
    \item Let's apply this to the Hopf map $h$. The first suspension is $E(h): S^4 \to S^3$. This gives an element in the group $\pi_4(S^3)$.
    \item It is a non-trivial result that $\pi_4(S^3) \cong \mathbb{Z}_2$. It is also a non-trivial result that the suspension homomorphism $E: \pi_3(S^2) \to \pi_4(S^3)$ is a map from $\mathbb{Z}$ to $\mathbb{Z}_2$ which sends the generator $h$ to the generator of $\mathbb{Z}_2$.
    \item This means that while you can wrap $S^3$ around $S^2$ an infinite number of distinct ways (classified by $\mathbb{Z}$), once you suspend the situation to mapping $S^4$ to $S^3$, all the even-integer wrappings become trivial, and all the odd-integer wrappings become equivalent. The stable phenomenon is the $\mathbb{Z}_2$ structure.
\end{itemize}

\textbf{3. The Freudenthal Suspension Theorem and Stable Stems}

The \textbf{Freudenthal Suspension Theorem} makes this precise. It says $E: \pi_{k+i}(S^k) \to \pi_{k+i+1}(S^{k+1})$ is an isomorphism when $k$ is large enough compared to $i$ (specifically, $k > i+1$). The groups $\pi_{k+i}(S^k)$ eventually become constant for a fixed "stem" $i = n-k$. This stable limit is the \textbf{$i$-th stable homotopy group of spheres}, $\pi_i^s$.
\begin{itemize}
    \item The stem $i=1$ stabilizes to $\pi_1^s = \pi_{k+1}(S^k)$ for $k \ge 2$. We have $\pi_3(S^2)=\mathbb{Z}$, but $\pi_4(S^3)=\mathbb{Z}_2$, and $\pi_5(S^4)=\mathbb{Z}_2$. The stable group is $\pi_1^s = \mathbb{Z}_2$.
\end{itemize}

\textbf{4. The Impossibility of a General Formula}

A full "classification" of $\pi_n(S^k)$ for all $n,k$ is a holy grail of topology; it is not known and not expected to have a simple answer. To illustrate the complexity, here is a table of the first few unstable groups.

\begin{center}
\begin{tabular}{c|ccccccccc}
$k\backslash n$ & 1 & 2 & 3 & 4 & 5 & 6 & 7 & 8 & 9 \\
\hline
1 & $\mathbb{Z}$ & 0 & 0 & 0 & 0 & 0 & 0 & 0 & 0 \\
2 & 0 & $\mathbb{Z}$ & $\mathbb{Z}$ & $\mathbb{Z}_2$ & $\mathbb{Z}_2$ & $\mathbb{Z}_{12}$ & $\mathbb{Z}_2$ & $\mathbb{Z}_2$ & $\mathbb{Z}_3$ \\
3 & 0 & 0 & $\mathbb{Z}$ & $\mathbb{Z}_2$ & $\mathbb{Z}_2$ & $\mathbb{Z}_{12}$ & $\mathbb{Z}_2$ & $\mathbb{Z}_2$ & $\mathbb{Z}_3$ \\
4 & 0 & 0 & 0 & $\mathbb{Z}$ & $\mathbb{Z}_2$ & $\mathbb{Z}_2$ & $\mathbb{Z}_{12}$ & $\mathbb{Z}_2$ & $\mathbb{Z}_2 \oplus \mathbb{Z}_2$ \\
5 & 0 & 0 & 0 & 0 & $\mathbb{Z}$ & $\mathbb{Z}_2$ & $\mathbb{Z}_2$ & $\mathbb{Z}_{24}$ & $\mathbb{Z}_2$ \\
6 & 0 & 0 & 0 & 0 & 0 & $\mathbb{Z}$ & $\mathbb{Z}_2$ & $\mathbb{Z}_2$ & $\mathbb{Z}_{24}$ \\
\end{tabular}
\end{center}
Notice how the groups along the anti-diagonals (constant $n-k=i$) eventually stabilize. For example, the stem $i=1$ reads $\pi_2(S^1)=0$, $\pi_3(S^2)=\mathbb{Z}$, $\pi_4(S^3)=\mathbb{Z}_2$, $\pi_5(S^4)=\mathbb{Z}_2$, etc., stabilizing at $\mathbb{Z}_2$.

\textbf{5. Further Structure and Computational Tools}
While there is no simple formula, there are deep structural results and powerful computational tools.
\begin{itemize}
    \item The groups $\pi_n(S^k)$ are all finite for $n>k$, except for groups of the form $\pi_{4m-1}(S^{2m})$ which contain a $\mathbb{Z}$ summand. These are related to the \textbf{J-homomorphism}, which connects the homotopy groups of spheres to the homotopy groups of special orthogonal groups.
    \item The primary tool for computing these groups is the \textbf{Adams spectral sequence}, which provides an intricate calculational framework relating the homotopy groups to more computable algebraic invariants from homological algebra.
\end{itemize}
In summary, the computation of $\pi_n(S^k)$ is an ongoing endeavor that sits at the heart of modern algebraic topology, revealing deep and often surprising connections between different areas of mathematics.
\end{remark}

\clearpage

\begin{exercise}
\label{geometry:2012:2}
Let $M$ be a smooth manifold of dimension $n$, and $X_{1}, \ldots, X_{k}$ be $k$ everywhere linearly independent smooth vector fields on an open set $U \subset M$ satisfying that $[X_{i}, X_{j}] = 0$ for $1 \leq i, j \leq k$. Prove that for any point $p \in U$ there is a coordinate chart $(V, y^{i})$ with $p \in V \subseteq U$ such that $X_{i} = \frac{\partial}{\partial y^{i}}$ for each $1 \leq i \leq k$.
\end{exercise}

\begin{solution}
\textbf{Geometric Intuition:} This is a special case of the Frobenius Theorem.\sidenote{The \textbf{Frobenius Theorem} states that a smooth distribution $\mathcal{D}$ on a manifold is integrable if and only if it is involutive.
\begin{itemize}
    \item A $k$-dimensional \textbf{distribution} $\mathcal{D}$ is a smooth assignment of a $k$-dimensional subspace $\mathcal{D}_p \subset T_pM$ at each point $p$. In this problem, the distribution is $\mathcal{D}_p = \text{span}\{X_1(p), \dots, X_k(p)\}$.
    \item A distribution is \textbf{involutive} if for any two vector fields $X, Y$ that lie in the distribution, their Lie bracket $[X, Y]$ also lies in the distribution. The given condition $[X_i, X_j] = 0$ for all $i,j$ is a strong form of this. It immediately implies that for any two vector fields $X = \sum f_i X_i$ and $Y = \sum g_j X_j$ in $\mathcal{D}$, their bracket $[X,Y]$ is also in $\mathcal{D}$.
    \item A distribution is \textbf{integrable} if, for every point, there is a coordinate chart $(V, y^i)$ where the distribution is spanned by the first $k$ coordinate vector fields, i.e., $\mathcal{D} = \text{span}\{\frac{\partial}{\partial y^1}, \dots, \frac{\partial}{\partial y^k}\}$.
\end{itemize}
The theorem guarantees that the "infinitesimal" algebraic condition of involutivity is the necessary and sufficient condition for the "local" geometric existence of these special coordinate charts (integral submanifolds).} The condition $[X_i, X_j]=0$ means that the flows of the vector fields commute, which allows us to build a consistent "grid" or coordinate system where the vector fields $X_i$ serve as the coordinate axes.

\textbf{A More Readable Proof (by Induction):}

The proof proceeds by induction on $k$, the number of vector fields.

\textbf{Base Case (k=1):} We are given a single non-vanishing vector field $X_1$. The "Flowbox Theorem" (or "Straightening Theorem") states that for any non-vanishing smooth vector field, we can always find local coordinates $(z^1, \dots, z^n)$ around any point $p$ such that in these coordinates, the vector field is constant: $X_1 = \frac{\partial}{\partial z^1}$.\sidenote{\textbf{Flowbox Theorem Statement:} Let $X$ be a smooth vector field on an $n$-manifold $M$. If $X(p) \neq 0$ for some $p \in M$, then there exists a coordinate chart $(U, (y^1, \dots, y^n))$ centered at $p$ such that on $U$, the vector field $X$ is given by $X = \frac{\partial}{\partial y^1}$. Intuitively, this theorem says that the integral curves of any non-vanishing vector field can be locally "straightened out" to become parallel to the first coordinate axis.} This establishes our base case.

\textbf{Inductive Step:} Assume the theorem is true for any set of $k-1$ commuting, linearly independent vector fields. Now we prove it for $k$ fields, $\{X_1, \dots, X_k\}$.

\begin{enumerate}
    \item \textbf{Straighten the first vector field:} From the base case, we know there exists a coordinate chart $(U, z^i)$ around $p$ such that $X_1 = \frac{\partial}{\partial z^1}$.

    \item \textbf{Analyze the remaining fields in the new chart:} Let's express the other vector fields $X_j$ (for $j=2, \dots, k$) in this new coordinate system:
    \[ X_j = \sum_{l=1}^n f_{jl}(z) \frac{\partial}{\partial z^l} \]
    where $f_{jl}(z)$ are smooth functions of the coordinates $(z^1, \dots, z^n)$.

    \item \textbf{Use the commuting property:} We are given that $[X_1, X_j] = 0$. Let's compute this Lie bracket:
    \[\begin{aligned}
    [X_1, X_j]
    &= \left[ \frac{\partial}{\partial z^1}, \sum_{l=1}^n f_{jl}(z) \frac{\partial}{\partial z^l} \right] \\
    &= \sum_{l=1}^n \left( \frac{\partial}{\partial z^1} f_{jl}(z) \right) \frac{\partial}{\partial z^l}
       - \sum_{l=1}^n f_{jl}(z) \underbrace{\left[\frac{\partial}{\partial z^1}, \frac{\partial}{\partial z^l}\right]}_{=0} \\
    &= \sum_{l=1}^n \frac{\partial f_{jl}}{\partial z^1} \frac{\partial}{\partial z^l} = 0
    \end{aligned}\]
    Since we know $[X_1, X_j]=0$, every component of this resulting vector field must be zero. This means:
    \[ \frac{\partial f_{jl}}{\partial z^1} = 0 \quad \text{for all } j \in \{2, \dots, k\} \text{ and } l \in \{1, \dots, n\}. \]
    This is a crucial simplification: the component functions $f_{jl}$ of the vector fields $X_2, \dots, X_k$ do not depend on the first coordinate, $z^1$.

    \item \textbf{Reduce the dimension:} This means that on any hypersurface defined by $z^1 = \text{constant}$, the vector fields $X_2, \dots, X_k$ are well-defined and depend only on the coordinates $(z^2, \dots, z^n)$. Furthermore, since $\{X_1, \dots, X_k\}$ are linearly independent and $X_1$ is not in the span of $\{\frac{\partial}{\partial z^2}, \dots, \frac{\partial}{\partial z^n}\}$, the vector fields $\{X_2, \dots, X_k\}$ when restricted to this hypersurface must be linearly independent and still commute with each other.

    \item \textbf{Apply the inductive hypothesis:} We now have a set of $k-1$ commuting, linearly independent vector fields on an $(n-1)$-dimensional manifold (the slice $z^1=\text{const}$). By our inductive hypothesis, there exists a coordinate change from $(z^2, \dots, z^n)$ to a new set $(y^2, \dots, y^n)$ such that in these new coordinates, $X_j = \frac{\partial}{\partial y^j}$ for $j=2, \dots, k$.
    Importantly, this coordinate change does not involve $z^1$. So we can define a new coordinate system for our original $n$-dimensional manifold:
    \[ y^1 = z^1, \quad y^2 = y^2(z^2, \dots, z^n), \quad \dots, \quad y^n = y^n(z^2, \dots, z^n). \]
    \item \textbf{Conclusion:} In this final $(y^1, \dots, y^n)$ coordinate system, we have:
    \begin{itemize}
        \item $X_1 = \frac{\partial}{\partial z^1} = \frac{\partial}{\partial y^1}$ (since $y^1$ is just $z^1$ and the other $y^i$ don't depend on $z^1$).
        \item $X_j = \frac{\partial}{\partial y^j}$ for $j=2, \dots, k$ (by the inductive step).
    \end{itemize}
    We have successfully constructed a coordinate chart where $X_i = \frac{\partial}{\partial y^i}$ for all $i=1, \dots, k$. The induction is complete.
\end{enumerate}
\end{solution}

\clearpage

\begin{exercise}
\label{geometry:2012:3}
Show that any self homeomorphism of $\mathbb{CP}^{2}$ is orientation preserving.\sidenote{A map $f: M \to M$ of an oriented $n$-manifold is \textbf{orientation-preserving} if it preserves the \textbf{local orientation} at each point $p \in M$, which is a choice of generator for $H_n(M, M \setminus \{p\}; \mathbb{Z}) \cong \mathbb{Z}$. In terms of tangent spaces, it is an equivalence class of ordered bases for $T_pM$ (bases are equivalent if the change-of-basis matrix has positive determinant). For a homeomorphism, being orientation-preserving is equivalent to having \textbf{degree} $\deg(f) = +1$. Conversely, a map is \textbf{orientation-reversing} if $\deg(f) = -1$. For complex manifolds, the natural orientation is induced by the complex structure $J$; a basis $\{v_1, Jv_1, \dots, v_n, Jv_n\}$ is always positively oriented.}
\end{exercise}

\begin{solution}
\textbf{Geometric Intuition:} Complex projective spaces carry a natural orientation from their complex structure.\sidenote{The $n$-dimensional complex projective space, $\mathbb{CP}^n$, is the space of all complex lines through the origin in $\mathbb{C}^{n+1}$.
\begin{itemize}
    \item \textbf{Quotient Construction:} It is formally defined as the quotient space $(\mathbb{C}^{n+1} \setminus \{0\}) / \sim$, where two points $z, w \in \mathbb{C}^{n+1}\setminus\{0\}$ are equivalent if $z = \lambda w$ for some non-zero complex number $\lambda \in \mathbb{C}^*$.
    \item \textbf{Homogeneous Coordinates:} A point is denoted by an equivalence class $[z_0:z_1:\dots:z_n]$.
    \item \textbf{Manifold Structure:} $\mathbb{CP}^n$ is a compact complex manifold of complex dimension $n$ (and thus real dimension $2n$). A key property is that its complex structure naturally induces a canonical orientation.
\end{itemize}} For $\mathbb{CP}^n$, any homeomorphism must respect the intersection product. In the case of $\mathbb{CP}^2$, the intersection of two "lines" (copies of $\mathbb{CP}^1$) is a single point with a positive sign. Even if a map reverses the sign of the 2-dimensional generator, the squared generator (representing the orientation) remains positive.

\textbf{Proof:} The cohomology ring of $\mathbb{CP}^2$ is given by $H^*(\mathbb{CP}^2; \mathbb{Z}) \cong \mathbb{Z}[x] / (x^3)$, where $x \in H^2(\mathbb{CP}^2; \mathbb{Z})$ is a generator. The orientation of $\mathbb{CP}^2$ is determined by the fundamental class $[M] \in H_4(\mathbb{CP}^2; \mathbb{Z})$, which is dual to the generator $x^2 \in H^4(\mathbb{CP}^2; \mathbb{Z})$.\footnote{
\textbf{Orientation and the Fundamental Class:}
For any compact, connected, $m$-dimensional manifold $M$, its top homology group $H_m(M; \mathbb{Z})$ is isomorphic to $\mathbb{Z}$.
\begin{enumerate}
    \item An \textbf{orientation} on $M$ is equivalent to choosing one of the two generators of this group. This chosen generator is called the \textbf{fundamental class}, denoted $[M]$.
    \item For $\mathbb{CP}^2$ (a real 4-manifold), this means choosing a generator for $H_4(\mathbb{CP}^2; \mathbb{Z}) \cong \mathbb{Z}$.
\end{itemize}
\textbf{Duality with Cohomology:}
There's a natural evaluation pairing (the Kronecker pairing) between cohomology and homology, $\langle \cdot, \cdot \rangle: H^k(M) \times H_k(M) \to \mathbb{Z}$.
\begin{enumerate}
    \item Once we've chosen a fundamental class $[M] \in H_4(\mathbb{CP}^2)$, it singles out a preferred generator for the top cohomology group $H^4(\mathbb{CP}^2; \mathbb{Z})$. This preferred generator, let's call it $[M]^\vee$, is the one that evaluates to 1 on the fundamental class: $\langle [M]^\vee, [M] \rangle = 1$. This is what "dual" means in this context.
    \item The cohomology ring structure provides a canonical generator for $H^4(\mathbb{CP}^2; \mathbb{Z})$, namely $x^2$, where $x$ is the generator of $H^2$ (dual to a projective line $\mathbb{CP}^1 \subset \mathbb{CP}^2$).
    \item We therefore \textit{define} the canonical orientation of $\mathbb{CP}^2$ by choosing the fundamental class $[M]$ to be the one that pairs with $x^2$ to give $+1$. This standard orientation is the one inherited from the complex structure.
\end{enumerate}
}

Let $f: \mathbb{CP}^2 \to \mathbb{CP}^2$ be a homeomorphism. It induces an automorphism $f^*: H^*(\mathbb{CP}^2; \mathbb{Z}) \to H^*(\mathbb{CP}^2; \mathbb{Z})$. Since $f^*$ must preserve the ring structure and $H^2(\mathbb{CP}^2; \mathbb{Z}) \cong \mathbb{Z}$, we must have $f^*(x) = \epsilon x$, where $\epsilon = \pm 1$.
Then for the generator of $H^4$, we have:
\[ f^*(x^2) = (f^*(x))^2 = (\epsilon x)^2 = \epsilon^2 x^2 = x^2. \]
The degree of $f$, denoted $\deg(f)$, is the integer such that $f^*([M]^\vee) = \deg(f) [M]^\vee$, where $[M]^\vee$ is the cohomological fundamental class (the generator of $H^4$). Our calculation shows $f^*(x^2) = 1 \cdot x^2$, so $\deg(f) = 1$. A map with degree $1$ is orientation preserving. Thus, any self homeomorphism of $\mathbb{CP}^2$ is orientation preserving.
\end{solution}
\begin{remark}[Intuitive Derivation of $H_*(\mathbb{CP}^n)$ and $H^*(\mathbb{CP}^n)$]
To understand why $\mathbb{CP}^n$ has this specific structure, we can look at its geometric construction and how it "builds up" from lower dimensions.

\textbf{1. Homology via Cell Decomposition (CW Complex):}
The most intuitive way to "see" the homology is to build $\mathbb{CP}^n$ inductively.
\begin{itemize}
    \item Start with a point: $\mathbb{CP}^0 = \{*\}$. This is a 0-cell ($e^0$).
    \item To get $\mathbb{CP}^1$ (the Riemann sphere $S^2$), we attach a 2-cell ($e^2 \cong \mathbb{C}$) to $\mathbb{CP}^0$.
    \item In general, $\mathbb{CP}^n$ is obtained by attaching a $2n$-cell ($e^{2n} \cong \mathbb{C}^n$) to $\mathbb{CP}^{n-1}$.
\end{itemize}
This gives a CW complex structure with exactly one cell in each even dimension: $e^0, e^2, \dots, e^{2n}$. Since there are no cells in odd dimensions, the cellular boundary maps $d_k$ must all be zero (as they map to or from a zero group). Thus, the homology is simply $\mathbb{Z}$ in every even dimension $0 \le 2k \le 2n$ and $0$ otherwise.

\textbf{2. Cohomology and the Ring Structure via Intersection:}
While homology tells us what "shapes" exist, cohomology (and the cup product) tells us how they \textbf{intersect}.
\begin{itemize}
    \item Let $x \in H^2(\mathbb{CP}^n)$ be the class dual to a hyperplane $L \cong \mathbb{CP}^{n-1} \subset \mathbb{CP}^n$.
    \item The cup product $x \cup x = x^2$ corresponds to the geometric intersection of two such hyperplanes in general position. In $\mathbb{CP}^n$, two hyperplanes intersect in a linear subspace of one lower dimension: $\mathbb{CP}^{n-1} \cap \mathbb{CP}^{n-1} \cong \mathbb{CP}^{n-2}$.
    \item Repeating this, $x^k$ corresponds to the intersection of $k$ hyperplanes, which is a $\mathbb{CP}^{n-k}$.
    \item Finally, $x^n$ is the intersection of $n$ hyperplanes, which is a single point (the fundamental class). Any further intersection $x^{n+1}$ is empty, hence $x^{n+1} = 0$.
\end{itemize}
This "intersection theory" motivation is why $H^*(\mathbb{CP}^n) \cong \mathbb{Z}[x]/(x^{n+1})$. In Exercise 3, we used $n=2$, so $x^2$ is the "point" class (the orientation), and the fact that $(-1)^2 = 1$ simply means that reflecting a line doesn't change the fact that two lines still intersect at a single positive point.
\textbf{Summary Comparison:}
\begin{center}
\begin{tabular}{|l|p{4cm}|p{5cm}|}
\hline
\textbf{Feature} & \textbf{Homology ($H_n$)} & \textbf{Cohomology ($H^n$)} \\ \hline
\textbf{Intuition} & Cycles/Chains (geometric objects) & Cocycles (functions/integrands) \\ \hline
\textbf{Orientation} & Defines the Fundamental Class $[M]$ & Dual to $[M]$ via Poincaré Duality \\ \hline
\textbf{Maps} & Push-forward ($f_*$) & Pull-back ($f^*$) \\ \hline
\textbf{Structure} & Group only & \textbf{Ring} (can "multiply" classes) \\ \hline
\end{tabular}
\end{center}

The crucial point for Exercise 3 is the \textbf{pull-back} $f^*$ and the \textbf{ring structure}. Because $f$ is a homeomorphism, it must induce an automorphism of the cohomology \textit{ring}. The fact that $H^4$ is generated by the square of a generator of $H^2$ allows the calculation $(-1)^2 = 1$, proving the map is orientation-preserving even if it reflects the 2-dimensional class.
\end{remark}

\clearpage

\begin{exercise}
\label{geometry:2012:4}
Prove the following version of the isoperimetric inequality: Suppose $C$ is a simple (that is, without self-intersection), smooth, closed curve in the Euclidean plane, with length $L$. Show that the area enclosed by $C$ is less than or equal to $\frac{L^{2}}{4\pi}$, and the equality occurs when and only if $C$ is a round circle.
\end{exercise}

\begin{solution}
\textbf{Geometric Intuition:} The circle is the "most efficient" shape in the sense that it maximizes area for a fixed perimeter. Any deviation from a circle (like an oval or a star) involves "stretching" the perimeter more than is necessary to contain the area.

\textbf{Proof:} Assume the curve $C$ is parametrized by arc length $s \in [0, L]$, so $(x(s), y(s))$ satisfies $(x')^2 + (y')^2 = 1$. The area $A$ is given by Green's Theorem: $A = \int_0^L x y' ds$. Without loss of generality, assume $\int_0^L x ds = 0$.
By Wirtinger's Inequality, if $f$ is $L$-periodic and $\int_0^L f ds = 0$, then $\int_0^L (f')^2 ds \geq (2\pi/L)^2 \int_0^L f^2 ds$.\sidenote{
\textbf{Derivation of Wirtinger's Inequality:}
Let $f(s)$ be an $L$-periodic function such that $\int_0^L f(s) ds = 0$. We can represent $f(s)$ using its Fourier series:
\[ f(s) = \sum_{n \in \mathbb{Z}} c_n e^{i \frac{2\pi n s}{L}} \]
where the condition $\int_0^L f(s) ds = 0$ implies that the constant term $c_0 = 0$.
The derivative is:
\[ f'(s) = \sum_{n \neq 0} c_n \left( i \frac{2\pi n}{L} \right) e^{i \frac{2\pi n s}{L}} \]
Applying \textbf{Parseval's Identity}, which relates the integral of the square of a function to the sum of the squares of its Fourier coefficients:
\[ \int_0^L f(s)^2 ds = L \sum_{n \neq 0} |c_n|^2 \]
\[ \int_0^L f'(s)^2 ds = L \sum_{n \neq 0} |c_n|^2 \left( \frac{2\pi n}{L} \right)^2 \]
Since $n^2 \geq 1$ for all $n \neq 0$ in the sum, we have:
\[ \int_0^L f'(s)^2 ds = \frac{4\pi^2}{L^2} L \sum_{n \neq 0} n^2 |c_n|^2 \geq \frac{4\pi^2}{L^2} L \sum_{n \neq 0} |c_n|^2 = \frac{4\pi^2}{L^2} \int_0^L f(s)^2 ds \]
This completes the derivation. Equality holds if and only if $c_n = 0$ for all $|n| > 1$, which means $f(s)$ is a pure sine or cosine wave of period $L$, i.e., $f(s) = a \cos(2\pi s/L) + b \sin(2\pi s/L)$.
}
We have $L = \int_0^L ((x')^2 + (y')^2) ds$.
Consider the quantity $L^2 - 4\pi A$:
\[ L^2 - 4\pi A = L \int_0^L ((x')^2 + (y')^2) ds - 4\pi \int_0^L x y' ds \]
Using the inequality $2xy' \leq \frac{2\pi}{L} x^2 + \frac{L}{2\pi} (y')^2$, we have:
\[ 4\pi A = 4\pi \int_0^L x y' ds \leq \int_0^L \left( \frac{4\pi^2}{L} x^2 + L (y')^2 \right) ds \]
Then:
\[ L^2 - 4\pi A \geq L \int_0^L (x')^2 ds + L \int_0^L (y')^2 ds - \int_0^L \frac{4\pi^2}{L} x^2 ds - L \int_0^L (y')^2 ds \]
\[ = L \int_0^L \left( (x')^2 - \left(\frac{2\pi}{L}\right)^2 x^2 \right) ds \geq 0 \]
by Wirtinger's inequality. Equality holds if and only if $x(s) = a \cos(2\pi s/L + \phi)$ and $y' = \frac{2\pi}{L} x$, which implies $y(s) = a \sin(2\pi s/L + \phi) + b$. This is the parametrization of a circle of radius $R = L/2\pi$.
\end{solution}

\clearpage

\begin{exercise}
\label{geometry:2012:5}
Let $x : M \to \mathbb{R}^{3}$ be a closed surface in 3-dimensional Euclidean space. Its Gaussian curvature and mean curvature are denoted by $K$ and $H$ respectively. Prove that:
\[ \iint_{M} H\, dA + \iint_{M} pK\, dA = 0, \quad \iint_{M} pH\, dA + \iint_{M} dA = 0, \]
where $p = \vec{x} \cdot \vec{n}$ is the support function of $M$, $\vec{x}$ denotes the position vector of $M$, $\vec{n}$ denotes the unit normal to $M$, and $dA$ is the area element of $M$.
\end{exercise}

\begin{solution}
\textbf{Geometric Intuition:} These formulas, known as the \textbf{Minkowski formulas}, reveal how a surface's "average" curvature is linked to its physical size and its distance from the origin.\footnote{\textbf{The Origin and Utility of Minkowski Formulas:} Hermann Minkowski originally developed these formulas to study "convex bodies." In modern geometry, they are used to prove "uniqueness" theorems—for example, showing that the only closed surface with constant mean curvature is a sphere. The formulas essentially say that the way a surface "bends" (curvature) must be perfectly balanced by its "spread" (area) and its "reach" (support function).}
\begin{itemize}
    \item The first formula, $\iint (H + pK) dA = 0$, relates the mean curvature $H$ and Gaussian curvature $K$ to the support function $p$.
    \item The second formula, $\iint (1 + pH) dA = 0$, is particularly intuitive: it tells us that on average, the support function $p$ and the mean curvature $H$ are inversely related.
\end{itemize}

\textbf{Preliminaries: Defining our Tools}
To prove these, we need to precisely define several geometric objects on our surface $M$:
\begin{enumerate}
    \item \textbf{Position Vector ($\vec{x}$):} A vector pointing from the origin $(0,0,0)$ to a point on the surface.
    \item \textbf{Unit Normal ($\vec{n}$):} A vector of length 1 that is perpendicular to the surface at each point.
    \item \textbf{Support Function ($p$):} Defined as $p = \vec{x} \cdot \vec{n}$. It represents the perpendicular distance from the origin to the tangent plane at $\vec{x}$.
    \item \textbf{Shape Operator ($S$):} This operator measures how the normal vector $\vec{n}$ changes as we move across the surface. It is defined by $S(v) = -\nabla_v \vec{n}$ for any tangent vector $v$.
    \item \textbf{Curvatures ($H$ and $K$):} $H = \frac{1}{2}\text{tr}(S)$ is the average of the "bending" rates, and $K = \det(S)$ is the product of the "bending" rates.
\end{enumerate}

\textbf{Step 1: Proving the Second Formula $\iint (1+pH) dA = 0$}
We use a common trick in global geometry: construct a vector field on the surface and apply the \textbf{Divergence Theorem}.\footnote{\textbf{The Divergence Theorem on Surfaces:} For any smooth vector field $X$ on a \textit{closed} (compact, no boundary) surface $M$, the total divergence must be zero: $\iint_M div(X) dA = 0$. This is the 2D version of the fact that if you have a flow on a balloon, the total "outflow" from all points must cancel out because there is no "edge" for the flow to escape from.}

Consider the \textbf{tangential position vector field} $X$:
\[ X = \vec{x}_{tan} = \vec{x} - (\vec{x} \cdot \vec{n})\vec{n} = \vec{x} - p\vec{n} \]
This field $X$ represents the part of the position vector that lies "flat" against the surface. We calculate its divergence, $div(X) = 2 + 2pH$.\sidenote{
\textbf{Derivation of $div(X) = 2 + 2pH$:}
Let $\{e_1, e_2\}$ be an orthonormal basis for the tangent space. Then $div(X) = \sum_{i=1}^2 \langle \nabla_{e_i} X, e_i \rangle$.
Substituting $X = \vec{x} - p\vec{n}$:
\[ \begin{aligned}
div(X) &= \sum_i \langle \nabla_{e_i} (\vec{x} - p\vec{n}), e_i \rangle \\
&= \sum_i \langle \nabla_{e_i} \vec{x}, e_i \rangle - \sum_i \langle \nabla_{e_i} (p\vec{n}), e_i \rangle
\end{aligned} \]
1. Since $\vec{x}$ is the position vector in $\mathbb{R}^3$, its derivative in the direction $e_i$ is simply $e_i$. Thus $\sum_i \langle e_i, e_i \rangle = 1 + 1 = 2$.
2. For the second term, use the product rule: $\nabla_{e_i} (p\vec{n}) = (e_i p)\vec{n} + p \nabla_{e_i} \vec{n}$.
\[ \langle (e_i p)\vec{n} + p \nabla_{e_i} \vec{n}, e_i \rangle = (e_i p)\underbrace{\langle \vec{n}, e_i \rangle}_{0} + p \langle \nabla_{e_i} \vec{n}, e_i \rangle \]
By definition of the shape operator, $\nabla_{e_i} \vec{n} = -S(e_i)$.
Thus, $\langle \nabla_{e_i} \vec{n}, e_i \rangle = -\langle S(e_i), e_i \rangle = -h_{ii}$.
Summing up:
\[ div(X) = 2 - \sum_i p(-h_{ii}) = 2 + p(h_{11} + h_{22}) = 2 + 2pH \]
}
Integrating over the closed surface $M$:\footnote{The \textbf{Divergence Theorem} on a manifold states that $\int_M div(X) dA = \int_{\partial M} \langle X, \nu \rangle ds$, where $\nu$ is the outward unit normal to the boundary. For a \textbf{closed manifold}, the boundary $\partial M$ is empty ($\emptyset$). Intuitively, if you have a "flow" (the vector field $X$) on a surface with no edges, any "expansion" in one area must be balanced by "contraction" elsewhere, as there is no boundary for the flow to escape or enter. Thus, the total net divergence over the entire surface must be exactly zero.}
\[ \iint_M (2 + 2pH) dA = 0 \implies \iint_M (1 + pH) dA = 0 \]

\textbf{Step 2: Proving the First Formula $\iint (H + pK) dA = 0$}
To incorporate the Gaussian curvature $K$, we modify our vector field using the \textbf{adjoint shape operator} $\tilde{S}$.\footnote{\textbf{The Adjoint Operator $\tilde{S}$:} In 2D, if $S$ is a matrix $\begin{pmatrix} a & b \\ b & d \end{pmatrix}$, then $\tilde{S} = \begin{pmatrix} d & -b \\ -b & a \end{pmatrix}$. Note that $\text{tr}(\tilde{S}) = a+d = \text{tr}(S) = 2H$, and $\tilde{S} S = \det(S) I = K I$. This operator is crucial because it "rotates" the curvature information to simplify the divergence calculation.}

Define a new vector field $Y = \tilde{S} X$. We calculate its divergence, $div(Y) = 2H + 2pK$.\sidenote{
\textbf{Derivation of $div(Y) = 2H + 2pK$:}
The divergence of a tensor $T$ acting on a vector $X$ is $div(TX) = (div \, T) \cdot X + \text{tr}(T \circ \nabla X)$.
1. A deep result in surface theory (the \textbf{Mainardi-Codazzi equations}) implies that for a surface in $\mathbb{R}^3$, the adjoint shape operator is divergence-free: $div \tilde{S} = 0$.
2. We compute the trace term: $\text{tr}(\tilde{S} \circ \nabla X) = \sum_i \langle \tilde{S}(\nabla_{e_i} X), e_i \rangle$.
From our previous step, $\nabla_{e_i} X = e_i - p \nabla_{e_i} \vec{n} = e_i + p S(e_i)$ (ignoring the normal component which disappears in the inner product).
\[ \begin{aligned}
\sum_i \langle \tilde{S}(e_i + p S(e_i)), e_i \rangle &= \sum_i \langle \tilde{S}e_i, e_i \rangle + p \sum_i \langle \tilde{S}S e_i, e_i \rangle \\
&= \text{tr}(\tilde{S}) + p \, \text{tr}(\tilde{S} S)
\end{aligned} \]
Since $\text{tr}(\tilde{S}) = 2H$ and $\tilde{S}S = KI$:
\[ div(Y) = 2H + p \, \text{tr}(KI) = 2H + 2pK \]
}
Applying the Divergence Theorem to $Y$:
\[ \iint_M (2H + 2pK) dA = 0 \implies \iint_M (H + pK) dA = 0 \]
This completes the proof of both Minkowski formulas.
\end{solution}

\clearpage

\begin{exercise}
\label{geometry:2012:6}
Write the structure equation of an orthonormal frame on a Riemannian manifold. Prove the following Riemannian metric $g$ has constant sectional curvature\sidenote{For a 2-dimensional subspace (a plane) $\sigma \subset T_pM$, the \textbf{sectional curvature} $K(\sigma)$ is defined as the Gaussian curvature of the surface $S = \exp_p(U)$ where $U$ is a neighborhood of the origin in $\sigma$ and $\exp_p$ is the exponential map. In terms of the Riemann curvature tensor $R$, if $\{u, v\}$ is an orthonormal basis for $\sigma$, then $K(\sigma) = R(u, v, v, u)$. A manifold has \textbf{constant sectional curvature} $c$ if $K(\sigma) = c$ for all planes $\sigma$ at all points $p$.} $c$ using the structure equation:
\[ g = \frac{\sum_{i=1}^{n} (dx^{i})^{2}}{[1 + \frac{c}{4}\sum_{i=1}^{n} (x^{i})^{2}]^{2}} \]
where $(x^{1}, \ldots, x^{n})$ is a local coordinate system.
\end{exercise}

\begin{solution}
\textbf{Geometric Intuition:} This metric, often called the \textbf{Riemannian model of constant curvature}, is a generalization of how we describe curved spaces like spheres or hyperboloids in flat coordinates.\footnote{\textbf{The Conformally Flat Metric:} A metric of the form $g = \phi^{-2} \sum (dx^i)^2$ is called "conformally flat" because it is a scaling of the flat Euclidean metric. The factor $\phi = 1 + \frac{c}{4}r^2$ is precisely what creates the curvature. If $c=0$, we have flat Euclidean space; if $c > 0$, we have an $n$-sphere (via stereographic projection); if $c < 0$, we have a hyperbolic space.} The goal is to show that no matter which "slice" or plane we look at, the surface "bends" with the same constant value $c$.

\textbf{1. The Structure Equations of Cartan}
To compute curvature without using messy Christoffel symbols, we use \textbf{moving frames}. We choose a set of "ruler" forms $\{\theta^i\}$ (a coframe) such that $g = \sum \theta^i \otimes \theta^i$.
\begin{enumerate}
    \item \textbf{First Structure Equation (Torsion-free):} $d\theta^i = \sum_j \omega^i_j \wedge \theta^j$. This defines the \textbf{connection 1-forms} $\omega^i_j$, which must be skew-symmetric ($\omega^i_j = -\omega^j_i$).
    \item \textbf{Second Structure Equation (Curvature):} $\Omega^i_j = d\omega^i_j - \sum_k \omega^i_k \wedge \omega^k_j$. These $\Omega^i_j$ are the \textbf{curvature 2-forms}.
\end{enumerate}
For a space with constant sectional curvature $c$, the curvature forms must satisfy the simple relationship: $\Omega^i_j = c \, \theta^i \wedge \theta^j$.

\textbf{2. Step-by-Step Proof}

\textbf{Step A: Defining the Coframe}
Given $g = \phi^{-2} \sum (dx^i)^2$, we define our orthonormal coframe as:
\[ \theta^i = \phi^{-1} dx^i, \quad \text{where } \phi = 1 + \frac{c}{4} \sum_{k=1}^n (x^k)^2. \]

\textbf{Step B: Finding the Connection 1-forms}
We need to find $\omega^i_j$ such that $d\theta^i = \sum \omega^i_j \wedge \theta^j$. Differentiating $\theta^i$ leads us to a specific form for $\omega^i_j$.\sidenote{
\textbf{Deriving the Connection Forms:}
1. Differentiate $\theta^i = \phi^{-1} dx^i$:
\[ d\theta^i = d(\phi^{-1}) \wedge dx^i = -\phi^{-2} d\phi \wedge dx^i = -\phi^{-1} d\phi \wedge \theta^i. \]
2. Calculate $d\phi$:
\[ d\phi = \frac{c}{4} d(\sum (x^k)^2) = \frac{c}{2} \sum x^k dx^k. \]
3. Substitute $dx^k = \phi \theta^k$:
\[ \phi^{-1} d\phi = \frac{c}{2} \sum x^k \theta^k. \]
4. Thus: $d\theta^i = -\frac{c}{2} (\sum_k x^k \theta^k) \wedge \theta^i = \sum_j \left( \frac{c}{2} (x^i \theta^j - x^j \theta^i) \right) \wedge \theta^j$.
5. By identifying the terms, we find: $\omega^i_j = \frac{c}{2} (x^i \theta^j - x^j \theta^i)$.
Notice this is skew-symmetric ($\omega^i_j = -\omega^j_i$), as required.
}

\textbf{Step C: Calculating the Curvature Forms}
Now we use the second structure equation $\Omega^i_j = d\omega^i_j - \sum_k \omega^i_k \wedge \omega^k_j$. This involves differentiating our $\omega^i_j$ and then subtracting the "wedge product" term.\sidenote{
\textbf{Algebraic Expansion of $\Omega^i_j$:}
1. \textbf{The $d\omega^i_j$ term:} Differentiate $\omega^i_j = \frac{c}{2} (x^i \theta^j - x^j \theta^i)$:
\[ d\omega^i_j = \frac{c}{2} (dx^i \wedge \theta^j + x^i d\theta^j - dx^j \wedge \theta^i - x^j d\theta^i). \]
Using $dx^i = \phi \theta^i$ and $d\theta^i = -\phi^{-1} d\phi \wedge \theta^i$:
\[ d\omega^i_j = c\phi \theta^i \wedge \theta^j - \frac{c}{2} \phi^{-1} d\phi \wedge (x^i \theta^j - x^j \theta^i). \]
2. \textbf{The $\omega \wedge \omega$ term:}
\[ \sum_k \omega^i_k \wedge \omega^k_j = \frac{c^2}{4} \sum_k (x^i \theta^k - x^k \theta^i) \wedge (x^k \theta^j - x^j \theta^k). \]
Expanding this product and using $\phi = 1 + \frac{c}{4} r^2$:
\[ \sum_k \omega^i_k \wedge \omega^k_j = \frac{c^2}{4} x^i (\sum x^k \theta^k) \wedge \theta^j + \frac{c^2}{4} x^j \theta^i \wedge (\sum x^k \theta^k) - \frac{c^2}{4} (\sum (x^k)^2) \theta^i \wedge \theta^j. \]
3. \textbf{Combining and Simplifying:}
When we subtract these terms, the terms involving $x^i, x^j$ cancel out perfectly, leaving:
\[ \Omega^i_j = (c\phi - \frac{c^2}{4} \sum (x^k)^2) \theta^i \wedge \theta^j. \]
Substituting $\phi = 1 + \frac{c}{4} \sum (x^k)^2$:
\[ \Omega^i_j = (c + \frac{c^2}{4} \sum (x^k)^2 - \frac{c^2}{4} \sum (x^k)^2) \theta^i \wedge \theta^j = c \, \theta^i \wedge \theta^j. \]
}

\textbf{Conclusion:}
Since the curvature forms satisfy $\Omega^i_j = c \, \theta^i \wedge \theta^j$, the sectional curvature of the metric $g$ is constant and equal to $c$.
\end{solution}

\clearpage

\subsection{Team Exam}

\begin{exercise}
\label{geometry:2012:7}
Prove that the real projective space $\mathbb{RP}^{n}$ is a differentiable manifold of dimension $n$.
\end{exercise}

\begin{solution}
\textbf{Geometric Intuition:} Real projective space $\mathbb{RP}^n$ is the space of all 1-dimensional linear subspaces (lines through the origin) in $\mathbb{R}^{n+1}$. A line is uniquely determined by any non-zero point $x$ on it, with the understanding that $x$ and $\lambda x$ represent the same line for any $\lambda \neq 0$. Alternatively, one can view $\mathbb{RP}^n$ as the $n$-sphere $S^n$ with antipodal points identified, $x \sim -x$. This identification "folds" the sphere such that each point in the resulting space represents a unique direction in the $(n+1)$-dimensional Euclidean space.

\textbf{1. Topological Manifold Structure}
We define $\mathbb{RP}^n$ as the quotient space $(\mathbb{R}^{n+1} \setminus \{0\}) / \sim$, where $x \sim y$ if $y = \lambda x$ for some $\lambda \in \mathbb{R} \setminus \{0\}$.\footnote{\textbf{The Quotient Topology and Projections:} For a topological space $X$ and an equivalence relation $\sim$, the quotient space $X/\sim$ is the set of equivalence classes equipped with the finest topology that makes the projection $\pi: X \to X/\sim$ continuous. Specifically, a set $V \subseteq X/\sim$ is defined to be open if and only if its preimage $\pi^{-1}(V)$ is an open set in the original space $X$. In our case, the projection $\pi$ is also an \textbf{open map}, meaning it maps open sets in $\mathbb{R}^{n+1} \setminus \{0\}$ to open sets in $\mathbb{RP}^n$.} A point in $\mathbb{RP}^n$ is represented by its \textbf{homogeneous coordinates} $[x_0 : x_1 : \dots : x_n]$.\footnote{\textbf{Homogeneous Coordinates and the Projective Line:} The notation $[x_0 : \dots : x_n]$ signifies that only the ratios between coordinates are geometrically significant. For example, in the projective line $\mathbb{RP}^1$, the points $[1:2]$ and $[2:4]$ are identical. This coordinate system is "homogeneous" because any homogeneous polynomial $P(x_0, \dots, x_n)$ has a well-defined zero set in $\mathbb{RP}^n$, even if the values of the polynomial depend on the choice of representative point.}

\textbf{Hausdorff and Second-Countable Properties:}
To show $\mathbb{RP}^n$ is a topological manifold, we check the following:
\begin{itemize}
    \item \textbf{Second-Countability:} The space $\mathbb{R}^{n+1} \setminus \{0\}$ is second-countable (it has a countable basis of open balls). Since the quotient map $\pi$ is an open map, it preserves the existence of a countable basis. Alternatively, consider the restriction $\pi|_{S^n}: S^n \to \mathbb{RP}^n$. Since $S^n$ is second-countable and $\pi$ is continuous and surjective, $\mathbb{RP}^n$ is second-countable.
    \item \textbf{Hausdorff Property:} $S^n$ is a compact Hausdorff space. The equivalence relation $x \sim \pm x$ corresponds to the action of the finite group $\mathbb{Z}_2 = \{Id, -Id\}$ on $S^n$ by homeomorphisms. The quotient of a Hausdorff space by a finite group action is always Hausdorff. Specifically, for any two distinct lines, we can find disjoint conical neighborhoods in $\mathbb{R}^{n+1} \setminus \{0\}$, which project to disjoint open sets in $\mathbb{RP}^n$.
\end{itemize}

\textbf{2. Differentiable Structure}
We must construct an atlas where the transition between any two charts is smooth.

\textbf{Construction of the Atlas:}
For each $i \in \{0, \dots, n\}$, define the subset:
\[ U_i = \{ [x_0 : \dots : x_n] \in \mathbb{RP}^n \mid x_i \neq 0 \}. \]
These sets $U_i$ are open because $\pi^{-1}(U_i) = \{ (x_0, \dots, x_n) \in \mathbb{R}^{n+1} \setminus \{0\} \mid x_i \neq 0 \}$ is open in $\mathbb{R}^{n+1}$. They cover $\mathbb{RP}^n$ because every non-zero vector must have at least one non-zero component.
Define the coordinate charts $\phi_i: U_i \to \mathbb{R}^n$ by:
\[ \phi_i([x_0 : \dots : x_n]) = \left( \frac{x_0}{x_i}, \dots, \frac{x_{i-1}}{x_i}, \frac{x_{i+1}}{x_i}, \dots, \frac{x_n}{x_i} \right). \]
This map is a homeomorphism. Its inverse $\phi_i^{-1}: \mathbb{R}^n \to U_i$ is given by:
\[ \phi_i^{-1}(y_1, \dots, y_n) = [y_1 : \dots : y_i : 1 : y_{i+1} : \dots : y_n] \]
where the entry $1$ is placed in the $i$-th position (index starting from $0$).

\textbf{Smoothness of Transition Maps:}
For the atlas to be smooth, the composition $\phi_j \circ \phi_i^{-1}$ must be a $C^\infty$ map for all $i, j$. Let $i < j$. The transition occurs on $U_i \cap U_j$, where both $x_i \neq 0$ and $x_j \neq 0$.\sidenote{
\textbf{Exact Derivation of $\phi_j \circ \phi_i^{-1}$:}
Let $y = (y_1, \dots, y_n) \in \phi_i(U_i \cap U_j)$. Then $y_j \neq 0$ because $x_j \neq 0$ and $x_i \neq 0$ (recalling $x_j$ corresponds to $y_j$ when $i < j$).
1. Apply the inverse chart:
$\phi_i^{-1}(y) = [y_1 : \dots : y_i : 1 : y_{i+1} : \dots : y_j : \dots : y_n]$.
2. Identify coordinates: $x_0 = y_1, \dots, x_i = 1, \dots, x_j = y_j, \dots, x_n = y_n$.
3. Apply the $j$-th chart:
$\phi_j([x]) = \left( \frac{x_0}{x_j}, \dots, \frac{x_{j-1}}{x_j}, \frac{x_{j+1}}{x_j}, \dots, \frac{x_n}{x_j} \right)$.
4. Substitute $y$ values:
$\phi_j(\phi_i^{-1}(y)) = \left( \frac{y_1}{y_j}, \dots, \frac{1}{y_j}, \dots, \frac{y_{j-1}}{y_j}, \frac{y_{j+1}}{y_j}, \dots, \frac{y_n}{y_j} \right)$.
Each component is a rational function of the coordinates $(y_1, \dots, y_n)$ with a non-vanishing denominator $y_j$ on the domain. Thus, the map is smooth.}
Since all transition maps are smooth, the collection $\{(U_i, \phi_i)\}_{i=0}^n$ forms a smooth atlas.

\textbf{Conclusion:}
$\mathbb{RP}^n$ is a Hausdorff, second-countable topological space equipped with a consistent smooth atlas of dimension $n$. Therefore, $\mathbb{RP}^n$ is a differentiable manifold of dimension $n$.
\end{solution}

\clearpage

\begin{exercise}
\label{geometry:2012:8}
Let $M, N$ be $n$-dimensional smooth, compact, connected manifolds, and $f : M \to N$ a smooth map with rank\footnote{The \textbf{rank} of a smooth map $f: M \to N$ at a point $x \in M$ is defined as the rank of the linear map $df_x: T_x M \to T_{f(x)} N$ between tangent spaces, which is the dimension of the image $\text{Im}(df_x)$. A map has \textbf{constant rank} if this value is the same for all $x \in M$. In this exercise, the rank is $n$ everywhere, which is the maximum possible rank (since $\dim M = \dim N = n$), making $f$ a \textbf{submersion} and an \textbf{immersion}, and thus a local diffeomorphism.} equals to $n$ everywhere. Show that $f$ is a covering map.\footnote{A surjective continuous map $p: E \to B$ is a \textbf{covering map} if every point $b \in B$ has an open neighborhood $V$ (called an \textbf{evenly covered neighborhood}) such that the preimage $p^{-1}(V)$ is a disjoint union of open sets $S_i \subset E$, each of which is mapped homeomorphically onto $V$ by $p$. These open sets $S_i$ are called the \textbf{sheets} of the covering over $V$. \textbf{Intuition:} If $V$ is a small disc, its preimage $p^{-1}(V)$ looks like a \textbf{stack of disjoint discs} (sheets) floating above it, where the map $p$ acts as a vertical projection that identifies each disc in the stack perfectly with the base disc $V$.}
\end{exercise}

\begin{solution}
\textbf{Geometric Intuition:} To understand why $f$ is a covering map, we first observe that the condition on the rank implies that $f$ is a local diffeomorphism.
\footnote{A smooth map $f: M \to N$ is a local diffeomorphism if every point $x \in M$ has an open neighborhood $U$ such that $f(U)$ is open in $N$ and the restriction $f|_U: U \to f(U)$ is a diffeomorphism.}
This means that locally, $M$ looks exactly like $N$, and $f$ simply identifies small patches of $M$ with patches of $N$. However, being a local diffeomorphism is not enough to be a covering map. For instance, the inclusion of an open interval into the circle is a local diffeomorphism but not a covering map because it doesn't "wrap around" properly.
The compactness of $M$ is the crucial ingredient that prevents the "edges" of $M$ from failing to cover $N$ or from accumulating in an infinite number of layers. Since $M$ is compact, the map must be "proper", meaning the preimage of any compact set is compact.
\footnote{A map $f: M \to N$ is said to be proper if the preimage $f^{-1}(K)$ of every compact set $K \subseteq N$ is compact in $M$. In the context of manifolds, a local diffeomorphism is a covering map if and only if it is a proper map.}
Combined with the connectedness of $N$, this ensures that $f$ wraps $M$ around $N$ a finite, constant number of times.

\textbf{Proof:} We proceed by verifying the conditions for a covering map in a step-by-step manner.

\textit{Step 1: $f$ is a local diffeomorphism.}
Since $\text{rank}(f) = n = \dim M = \dim N$ at every point $x \in M$, the derivative $df_x: T_x M \to T_{f(x)} N$ is a linear isomorphism between tangent spaces of the same dimension. By the Inverse Function Theorem, for every $x \in M$, there exists an open neighborhood $U_x$ such that $f|_{U_x}: U_x \to f(U_x)$ is a diffeomorphism.

\textit{Step 2: $f$ is surjective.}
First, we show that $f(M)$ is an open set. Since $f$ is a local diffeomorphism, for any $x \in M$, the image of its neighborhood $f(U_x)$ is open in $N$.\sidenote{A map $f: X \to Y$ is \textbf{open} if the image of every open set is open, and \textbf{closed} if the image of every closed set is closed. These properties are generally independent; a map can be one, both, or neither. In this problem, $f$ is \textbf{open} because it is a local diffeomorphism, and it is \textbf{closed} because it is a continuous map from a compact space to a Hausdorff space. Being both open and closed is what allows us to conclude $f(M) = N$ via connectedness.}
\sidenote{Since $f$ is locally a diffeomorphism, it is an open map; that is, the image of any open set in $M$ is open in $N$.}
Next, we show that $f(M)$ is a closed set.
\sidenote{A continuous map from a compact space to a Hausdorff space is always a closed map. Since $M$ is compact and $N$ is a manifold (and thus Hausdorff), $f$ maps closed sets to closed sets.}
Because $M$ is closed in itself, $f(M)$ is closed in $N$. Since $N$ is connected and $f(M)$ is a non-empty subset that is both open and closed, we must have $f(M) = N$.

\textit{Step 3: The fibers of $f$ are finite.}
For any $y \in N$, consider the fiber $f^{-1}(y)$. Since $f$ is continuous and $\{y\}$ is closed, $f^{-1}(y)$ is a closed subset of the compact manifold $M$, and is therefore compact.
\sidenote{In a local diffeomorphism, every point $x \in f^{-1}(y)$ has a neighborhood $U_x$ such that $f|_{U_x}$ is injective. This implies $U_x \cap f^{-1}(y) = \{x\}$, which shows the set is discrete.}
A compact set that is also discrete must be finite. Let $f^{-1}(y) = \{x_1, x_2, \dots, x_k\}$ for some $k \in \mathbb{N}$.

\textit{Step 4: Construction of an evenly covered neighborhood.}
We need to show that for any $y \in N$, there exists an open neighborhood $V$ that is evenly covered. Let $f^{-1}(y) = \{x_1, \dots, x_k\}$. By the local diffeomorphism property, we can choose disjoint open neighborhoods $U_1, \dots, U_k$ of $x_1, \dots, x_k$ respectively, such that each $f|_{U_i}: U_i \to V_i = f(U_i)$ is a diffeomorphism.
To ensure $V$ is evenly covered, we must exclude points in $N$ that are images of points in $M$ outside the union of $U_i$. Define $K = M \setminus \bigcup_{i=1}^k U_i$.
\sidenote{Since $\cup U_i$ is open, its complement $K$ is closed in $M$. As $M$ is compact, $K$ is compact, and its image $f(K)$ is closed in $N$.}
Note that $y \notin f(K)$ because all preimages of $y$ are contained in the $U_i$. Now, let
\[ V = \left( \bigcap_{i=1}^k V_i \right) \setminus f(K). \]
Since $V$ is the intersection of finitely many open sets $V_i$ minus a closed set $f(K)$, $V$ is an open neighborhood of $y$.
By construction, $f^{-1}(V) \cap K = \emptyset$, which means $f^{-1}(V) \subset \bigcup_{i=1}^k U_i$. Therefore,
\[ f^{-1}(V) = \bigcup_{i=1}^k (f^{-1}(V) \cap U_i). \]
Each set $W_i = f^{-1}(V) \cap U_i$ is an open neighborhood of $x_i$ that is mapped diffeomorphically onto $V$ by $f$. Since the $U_i$ were chosen to be disjoint, the $W_i$ are also disjoint. Thus, $V$ is an evenly covered neighborhood of $y$.

Since this holds for any $y \in N$ and $f$ is surjective, $f$ is a covering map.
\end{solution}

\clearpage

\begin{exercise}
\label{geometry:2012:9}
Given any Riemannian manifold $(M^{n}, g)$, show that there exists a unique Riemannian connection\sidenote{A connection $\nabla$ is the \textbf{Riemannian connection} (also called the \textbf{Levi-Civita connection}) if it is \textbf{metric-compatible} ($X g(Y, Z) = g(\nabla_X Y, Z) + g(Y, \nabla_X Z)$) and \textbf{torsion-free} ($\nabla_X Y - \nabla_Y X = [X, Y]$). It is a fundamental theorem of Riemannian geometry that such a connection exists and is unique for any Riemannian manifold.} on $M^{n}$.
\end{exercise}

\begin{solution}
\textbf{Geometric Intuition:} The problem of finding a "natural" way to differentiate vector fields on a manifold is central to Riemannian geometry. While there are infinitely many affine connections on any smooth manifold, the presence of a Riemannian metric $g$ allows us to distinguish a unique one. We seek a connection that is "compatible" with the metric and "straight" (torsion-free).
\begin{itemize}
    \item \textbf{Metric Compatibility:} This condition ensures that the inner product of two vector fields is preserved if they are parallel transported. Geometrically, this means the connection preserves lengths and angles, treating the metric as "constant" along any curve.
    \item \textbf{Torsion-free (Symmetry):} This condition ensures that the connection has no "intrinsic twisting." Geometrically, it means that "infinitesimal parallelograms" formed by the flows of vector fields close.
\end{itemize}
The Fundamental Theorem of Riemannian Geometry states that these two requirements uniquely characterize the Levi-Civita connection, making it the canonical derivative for Riemannian manifolds.

\textbf{Proof:} A Riemannian connection (or Levi-Civita connection) $\nabla$ on $(M, g)$ is an affine connection that is uniquely characterized by two fundamental properties: \textbf{metric compatibility} and the \textbf{torsion-free} condition.

\textbf{1. Uniqueness:}
Suppose such a connection $\nabla$ exists. To show it is unique, we demonstrate that its values are entirely determined by the metric $g$ and the Lie bracket. By cyclically permuting the variables in the metric compatibility condition, we obtain:
\begin{align*}
    (1) \quad Xg(Y, Z) &= g(\nabla_X Y, Z) + g(Y, \nabla_X Z) \\
    (2) \quad Yg(Z, X) &= g(\nabla_Y Z, X) + g(Z, \nabla_Y X) \\
    (3) \quad Zg(X, Y) &= g(\nabla_Z X, Y) + g(X, \nabla_Z Y)
\end{align*}
By considering the combination $(1) + (2) - (3)$ and substituting the torsion-free condition, we arrive at the \textbf{Koszul Formula}\sidenote{Derivation of the Koszul Formula:
Summing $(1) + (2) - (3)$ gives:
$Xg(Y,Z) + Yg(Z,X) - Zg(X,Y) = g(\nabla_X Y + \nabla_Y X, Z) + g(\nabla_X Z - \nabla_Z X, Y) + g(\nabla_Y Z - \nabla_Z Y, X)$
Using $\nabla_A B - \nabla_B A = [A, B]$:
$\nabla_X Y + \nabla_Y X = 2\nabla_X Y - [X, Y]$
$\nabla_X Z - \nabla_Z X = [X, Z]$
$\nabla_Y Z - \nabla_Z Y = [Y, Z]$
Substituting these into the sum:
$= g(2\nabla_X Y - [X, Y], Z) + g([X, Z], Y) + g([Y, Z], X)$
$= 2g(\nabla_X Y, Z) - g([X, Y], Z) + g([X, Z], Y) + g([Y, Z], X)$.
Rearranging for $2g(\nabla_X Y, Z)$ yields the result.}:
\[ \begin{aligned}
    2g(\nabla_X Y, Z) &= X g(Y, Z) + Y g(X, Z) - Z g(X, Y) \\
    &\quad + g([X, Y], Z) - g([X, Z], Y) - g([Y, Z], X).
\end{aligned} \]
Because the Riemannian metric $g$ is non-degenerate, the inner product $g(\nabla_X Y, Z)$ being specified for all $Z \in \mathfrak{X}(M)$ uniquely determines the vector field $\nabla_X Y$. Thus, the connection is unique.

\textbf{2. Existence:}
To prove existence, we define the operator $\nabla$ using the Koszul formula. For any fixed $X, Y \in \mathfrak{X}(M)$, the mapping $Z \mapsto \text{RHS}(X, Y, Z)$ is $C^\infty(M)$-linear. By the non-degeneracy of the Riemannian metric $g$, there exists a unique vector field, which we denote by $\nabla_X Y$, satisfying the Koszul formula for all $Z$:
\[ 2g(\nabla_X Y, Z) = X g(Y, Z) + Y g(X, Z) - Z g(X, Y) + g([X, Y], Z) - g([X, Z], Y) - g([Y, Z], X). \]
We must now verify that this implicitly defined $\nabla$ satisfies all the axioms of a Riemannian connection.

\textbf{Stage A: Verifying Connection Axioms}
A connection must be $C^\infty(M)$-linear in $X$ and satisfy the Leibniz rule in $Y$. We verify these properties by substituting the modified fields into the Koszul formula.

\begin{enumerate}
    \item \textbf{Linearity in $X$:} Replace $X$ with $fX$ in the formula for $2g(\nabla_{fX} Y, Z)$:
    \[ \begin{aligned}
    2g(\nabla_{fX} Y, Z) &= (fX)g(Y, Z) + Yg(Z, fX) - Zg(fX, Y) \\
    &\quad + g([fX, Y], Z) - g([fX, Z], Y) - g([Y, Z], fX)
    \end{aligned} \]
    Using the properties of the Lie bracket and the metric:
    \begin{itemize}
        \item $(fX)g(Y, Z) = f Xg(Y, Z)$
        \item $Yg(Z, fX) = Y(f)g(Z, X) + f Yg(Z, X)$
        \item $Zg(fX, Y) = Z(f)g(X, Y) + f Zg(X, Y)$
        \item $g([fX, Y], Z) = g(f[X, Y] - Y(f)X, Z) = f g([X, Y], Z) - Y(f)g(X, Z)$
        \item $g([fX, Z], Y) = g(f[X, Z] - Z(f)X, Y) = f g([X, Z], Y) - Z(f)g(X, Y)$
        \item $g([Y, Z], fX) = f g([Y, Z], X)$
    \end{itemize}
    Summing these terms:
    \[ \begin{aligned}
    2g(\nabla_{fX} Y, Z) &= f Xg(Y, Z) + Y(f)g(Z, X) + f Yg(Z, X) - Z(f)g(X, Y) - f Zg(X, Y) \\
    &\quad + f g([X, Y], Z) - Y(f)g(X, Z) - f g([X, Z], Y) + Z(f)g(X, Y) - f g([Y, Z], X)
    \end{aligned} \]
    Observe that $Y(f)g(Z, X) - Y(f)g(X, Z) = 0$ and $-Z(f)g(X, Y) + Z(f)g(X, Y) = 0$. The remaining terms are precisely $f \cdot 2g(\nabla_X Y, Z)$, which proves $C^\infty(M)$-linearity in $X$.

    \item \textbf{Leibniz Rule in $Y$:} We must show $\nabla_X(fY) = X(f)Y + f\nabla_X Y$. Substitute $fY$ into the formula for $2g(\nabla_X(fY), Z)$:
    \[ \begin{aligned}
    2g(\nabla_X(fY), Z) &= Xg(fY, Z) + fYg(Z, X) - Zg(X, fY) \\
    &\quad + g([X, fY], Z) - g([X, Z], fY) - g([fY, Z], X)
    \end{aligned} \]
    Expanding each term:
    \begin{itemize}
        \item $Xg(fY, Z) = X(f)g(Y, Z) + fXg(Y, Z)$
        \item $fYg(Z, X) = f Yg(Z, X)$
        \item $Zg(X, fY) = Z(f)g(X, Y) + f Zg(X, Y)$
        \item $g([X, fY], Z) = g(X(f)Y + f[X, Y], Z) = X(f)g(Y, Z) + f g([X, Y], Z)$
        \item $g([X, Z], fY) = f g([X, Z], Y)$
        \item $g([fY, Z], X) = g(f[Y, Z] - Z(f)Y, X) = f g([Y, Z], X) - Z(f)g(Y, X)$
    \end{itemize}
    Summing the expanded terms:
    \[ \begin{aligned}
    2g(\nabla_X(fY), Z) &= X(f)g(Y, Z) + f Xg(Y, Z) + f Yg(Z, X) - Z(f)g(X, Y) - f Zg(X, Y) \\
    &\quad + X(f)g(Y, Z) + f g([X, Y], Z) - f g([X, Z], Y) - f g([Y, Z], X) + Z(f)g(Y, X)
    \end{aligned} \]
    The $Z(f)$ terms cancel: $-Z(f)g(X, Y) + Z(f)g(Y, X) = 0$. The $X(f)$ terms combine: $X(f)g(Y, Z) + X(f)g(Y, Z) = 2X(f)g(Y, Z)$. The remaining terms are $f \cdot 2g(\nabla_X Y, Z)$. Thus, $2g(\nabla_X(fY), Z) = 2g(X(f)Y + f\nabla_X Y, Z)$, proving the Leibniz rule.
\end{enumerate}

\textbf{Stage B: Verifying Riemannian Properties}
We check that $\nabla$ satisfies the torsion-free property and is compatible with the metric.

\begin{enumerate}
    \item \textbf{Torsion-free Property ($\nabla_X Y - \nabla_Y X = [X, Y]$):} Subtract $2g(\nabla_Y X, Z)$ from $2g(\nabla_X Y, Z)$:
    \[ \begin{aligned}
    2g(\nabla_X Y - \nabla_Y X, Z) &= [Xg(Y, Z) + Yg(X, Z) - Zg(X, Y) + g([X, Y], Z) - g([X, Z], Y) - g([Y, Z], X)] \\
    &\quad - [Yg(X, Z) + Xg(Y, Z) - Zg(Y, X) + g([Y, X], Z) - g([Y, Z], X) - g([X, Z], Y)]
    \end{aligned} \]
    Performing the subtraction:
    \begin{itemize}
        \item $Xg(Y, Z) - Xg(Y, Z) = 0$
        \item $Yg(X, Z) - Yg(X, Z) = 0$
        \item $-Zg(X, Y) + Zg(Y, X) = 0$
        \item $g([X, Y], Z) - g([Y, X], Z) = 2g([X, Y], Z)$
        \item $-g([X, Z], Y) + g([X, Z], Y) = 0$
        \item $-g([Y, Z], X) + g([Y, Z], X) = 0$
    \end{itemize}
    We obtain $2g(\nabla_X Y - \nabla_Y X, Z) = 2g([X, Y], Z)$, which proves the torsion-free condition.

    \item \textbf{Metric Compatibility ($Xg(Y, Z) = g(\nabla_X Y, Z) + g(Y, \nabla_X Z)$):} Add $2g(\nabla_X Y, Z)$ and $2g(\nabla_X Z, Y)$:
    \[ \begin{aligned}
    2g(\nabla_X Y, Z) + 2g(\nabla_X Z, Y) &= [Xg(Y, Z) + Yg(X, Z) - Zg(X, Y) + g([X, Y], Z) - g([X, Z], Y) - g([Y, Z], X)] \\
    &\quad + [Xg(Z, Y) + Zg(X, Y) - Yg(X, Z) + g([X, Z], Y) - g([X, Y], Z) - g([Z, Y], X)]
    \end{aligned} \]
    Summing the terms:
    \begin{itemize}
        \item $Xg(Y, Z) + Xg(Z, Y) = 2Xg(Y, Z)$
        \item $Yg(X, Z) - Yg(X, Z) = 0$
        \item $-Zg(X, Y) + Zg(X, Y) = 0$
        \item $g([X, Y], Z) - g([X, Y], Z) = 0$
        \item $-g([X, Z], Y) + g([X, Z], Y) = 0$
        \item $-g([Y, Z], X) - g([Z, Y], X) = 0$ (by skew-symmetry of the bracket $[Y, Z] = -[Z, Y]$)
    \end{itemize}
    The sum is exactly $2Xg(Y, Z)$, which proves $g(\nabla_X Y, Z) + g(Y, \nabla_X Z) = Xg(Y, Z)$.
\end{enumerate}
Since the operator $\nabla$ defined by the Koszul formula satisfies all axioms of an affine connection, is torsion-free, and is compatible with the metric, existence is proved.
\end{solution}

\clearpage

\begin{exercise}
\label{geometry:2012:10}
Let $S^{n}$ be the unit sphere in $\mathbb{R}^{n+1}$ and $f : S^{n} \to S^{n}$ a continuous map. Assume that the degree of $f$ is an odd integer. Show that there exists $x_{0} \in S^{n}$ such that $f(-x_{0}) = -f(x_{0})$.
\end{exercise}

\begin{solution}
\textbf{Geometric Intuition:} This problem is a fundamental result in algebraic topology, closely related to the Borsuk-Ulam Theorem. The core idea is that a map of odd degree must "spread" across the sphere in such a way that it cannot avoid mapping at least one pair of antipodal points to antipodal points. If such a map were to avoid this condition, it could be continuously deformed (homotoped) into an even map.\sidenote{
\textbf{Definition of an Even Map:}
A map $g: S^n \to Y$ is called \textbf{even} if $g(x) = g(-x)$ for all $x \in S^n$. Geometrically, this means the map identifies every pair of antipodal points on the sphere. Consequently, such a map factors through the quotient space $\mathbb{RP}^n = S^n / \{x \sim -x\}$, where each point represents a pair of antipodal points.
} However, even maps factor through the real projective space $\mathbb{RP}^n$, which imposes a parity constraint: their degree must always be even.

\begin{remark}[Homotopy to Even Maps and Projective Factorization]
This proof relies on two deep topological insights:
\begin{enumerate}
    \item \textbf{Why we can deform $f$ into an even map:} The assumption for contradiction is that $f(x) \neq -f(-x)$ for all $x$. This means the vector sum $f(x) + f(-x)$ is never the zero vector. Because the sum is never zero, we can "slide" the value of $f(x)$ towards this average value along a straight line (and then project back to the sphere) without ever hitting the origin. The "midpoint" of this slide is the map $g(x) = \frac{f(x) + f(-x)}{|f(x) + f(-x)|}$. Since $f(x) + f(-x) = f(-x) + f(x)$, the map $g$ is perfectly symmetric ($g(x) = g(-x)$), making it an even map.
    \item \textbf{Why even maps factor through $\mathbb{RP}^n$:} By definition, the real projective space $\mathbb{RP}^n$ is the quotient of the sphere $S^n$ where each point $x$ is "glued" to its opposite point $-x$. A universal property of quotient spaces states that if a map $g: S^n \to Y$ treats equivalent points identically (i.e., $g(x) = g(-x)$), then $g$ doesn't "see" the difference between the sphere and the projective space. Thus, $g$ "factors" as $g = \bar{g} \circ p$, where $p: S^n \to \mathbb{RP}^n$ is the projection and $\bar{g}$ is a well-defined map on the projective space itself.
\end{enumerate}
\end{remark}

\textbf{Proof:} Assume for contradiction that there exists no $x_0 \in S^n$ such that $f(-x_0) = -f(x_0)$. This hypothesis implies that for all $x \in S^n$, the vectors $f(x)$ and $f(-x)$ are never antipodal in $\mathbb{R}^{n+1}$, which is analytically expressed as $f(x) \neq -f(-x)$, or equivalently, $f(x) + f(-x) \neq 0$ for all $x \in S^n$.

Since $f(x) + f(-x)$ is never the zero vector, the line segment connecting $f(x)$ and $f(-x)$ in $\mathbb{R}^{n+1}$ never passes through the origin. This allows us to define a continuous homotopy $H: S^n \times [0, 1] \to S^n$ between the map $f$ and its composition with the antipodal map\footnote{The antipodal map $A: S^n \to S^n$ is defined by $A(x) = -x$. It maps each point on the sphere to its diametrically opposite point. In the context of $S^n \subset \mathbb{R}^{n+1}$, $\deg(A) = (-1)^{n+1}$.} $A$ as follows:
\[ H(x, t) = \frac{(1-t)f(x) + t f(-x)}{|(1-t)f(x) + t f(-x)|} \]
At $t=0$, we have $H(x, 0) = f(x)$, and at $t=1$, we have $H(x, 1) = f(-x) = (f \circ A)(x)$. Thus, $f$ is homotopic to $f \circ A$, which implies that $\deg(f) = \deg(f \circ A)$.

\begin{definition}[Degree of a Map]
For a continuous map $f: S^n \to S^n$, the \textbf{degree} $\deg(f)$ is the unique integer $d$ such that the induced homomorphism on the top homology group:
\[ f_*: H_n(S^n; \mathbb{Z}) \to H_n(S^n; \mathbb{Z}) \]
satisfies $f_*(\alpha) = d\alpha$ for any class $\alpha$. Intuitively, the degree measures how many times the map "wraps" the sphere around itself. A key property is that \textbf{homotopic maps} always have equal degrees.
\end{definition}

\textbf{Calculation of $\deg(f \circ A)$:}
By the functoriality of the degree, we have $\deg(f \circ A) = \deg(f) \cdot \deg(A)$. Since $\deg(A) = (-1)^{n+1}$, it follows that $\deg(f \circ A) = (-1)^{n+1} \deg(f)$.
If $n$ is even, then $(-1)^{n+1} = -1$, which implies $\deg(f) = -\deg(f)$, hence $\deg(f) = 0$. Since $0$ is an even integer, this already contradicts the assumption that $\deg(f)$ is odd.

To handle the case where $n$ is odd, we note that at the midpoint of the homotopy ($t = 1/2$), we obtain a map $g: S^n \to S^n$ defined by:
\[ g(x) = \frac{f(x) + f(-x)}{|f(x) + f(-x)|} \]
This map $g$ is \textit{even}, meaning $g(x) = g(-x)$ for all $x \in S^n$. By construction, $f$ is homotopic to $g$, so $\deg(f) = \deg(g)$. We now show that any even map $g: S^n \to S^n$ must have an even degree.

\textbf{Degree of even maps:}
An even map $g$ factors as $g = \bar{g} \circ p$, where $p: S^n \to \mathbb{RP}^n$ is the standard covering map and $\bar{g}: \mathbb{RP}^n \to S^n$.
\begin{enumerate}
    \item If $n$ is even, $H_n(\mathbb{RP}^n; \mathbb{Z}) = 0$, so $g_*$ is the zero map and $\deg(g) = 0$.
    \item If $n$ is odd, $\mathbb{RP}^n$ is orientable. The map $p$ has degree $2$ as a $2$-sheeted cover. Thus:
    \[ \deg(g) = \deg(\bar{g}) \cdot \deg(p) = 2 \deg(\bar{g}). \]
\end{enumerate}
In both cases, $\deg(g)$ is an even integer.

Thus, the degree of $f$ must be even, which contradicts the given assumption that $\deg(f)$ is odd. Therefore, there must exist at least one point $x_0 \in S^n$ such that $f(-x_0) = -f(x_0)$.
\end{solution}

\clearpage

\begin{exercise}
\label{geometry:2012:11}
State and prove the Stokes theorem for oriented compact manifolds.
\end{exercise}

\begin{solution}
\textbf{Geometric Intuition:} Stokes' Theorem is the ultimate generalization of the Fundamental Theorem of Calculus. It says that the "total accumulation" of a derivative inside a region is perfectly balanced by the "flux" or value of the original function on the boundary.

\textbf{Statement of the Theorem:}
Let $M$ be an oriented compact smooth $n$-dimensional manifold with boundary $\partial M$, and let $\omega$ be a smooth $(n-1)$-form on $M$. Let $\partial M$ be given the induced boundary orientation.\footnote{
\textbf{Induced Boundary Orientation:}
For a point $p \in \partial M$, a basis $\{v_1, \dots, v_{n-1}\}$ for the tangent space $T_p(\partial M)$ is \textit{positively oriented} if the basis $\{n, v_1, \dots, v_{n-1}\}$ is a positively oriented basis for $T_p M$, where $n$ is an \textbf{outward-pointing} normal vector at $p$.
} Then:
\[ \int_M d\omega = \int_{\partial M} \omega. \]

\textbf{Proof:}
The proof proceeds by localizing the problem using a partition of unity and then verifying the equality in local coordinate charts.

\textbf{Stage 1: Localization via Partition of Unity}
Because $M$ is compact, we can cover it with a finite number of coordinate charts $\{U_\alpha, \phi_\alpha\}$. We employ a \textbf{partition of unity}\footnote{
\textbf{Partition of Unity:}
A partition of unity subordinate to a cover $\{U_\alpha\}$ is a collection of smooth functions $\{\rho_\alpha\}$ such that $0 \le \rho_\alpha \le 1$, the support of each $\rho_\alpha$ is contained in some $U_\alpha$, and $\sum \rho_\alpha = 1$ everywhere on $M$. This tool allows us to decompose a global form into a sum of locally supported forms: $\omega = \sum_\alpha (\rho_\alpha \omega)$.
} $\{\rho_\alpha\}$ subordinate to this cover.
By the linearity of the integral and the exterior derivative, it suffices to prove the theorem for a form $\eta = \rho_\alpha \omega$ that is compactly supported\footnote{
\textbf{Compact Support:}
A form $\eta$ has compact support if the closure of the set where $\eta \neq 0$ is compact. Since $M$ is compact, every smooth form on $M$ has compact support. Localization ensures that we only need to consider the behavior of $\eta$ within a single coordinate patch.
} in a single coordinate chart $U$.

\textbf{Stage 2: Case 1 - Support in an Interior Chart}
Suppose the support of $\eta$ is contained in an interior chart $U \cong \mathbb{R}^n$. We express $\eta$ in coordinates as:
\[ \eta = \sum_{i=1}^n (-1)^{i-1} a_i(x) dx_1 \wedge \dots \wedge \widehat{dx_i} \wedge \dots \wedge dx_n. \]
Taking the exterior derivative, we have $d\eta = \sum_{i=1}^n \frac{\partial a_i}{\partial x_i} dx_1 \wedge \dots \wedge dx_n$.
Since the support of $\eta$ does not reach the boundary, $\int_{\partial M} \eta = 0$.
\sidenote{
\textbf{Calculation for Interior Support:}
We compute the integral over $\mathbb{R}^n$:
\[ \int_{\mathbb{R}^n} d\eta = \sum_{i=1}^n \int_{\mathbb{R}^n} \frac{\partial a_i}{\partial x_i} dx_1 \dots dx_n \]
By Fubini's theorem, we can integrate each term $i$ with respect to $x_i$ first:
\[ \int_{\mathbb{R}^{n-1}} \left( \int_{-\infty}^{\infty} \frac{\partial a_i}{\partial x_i} dx_i \right) dx_1 \dots \widehat{dx_i} \dots dx_n \]
By the \textbf{Fundamental Theorem of Calculus}, the inner integral is zero because $a_i$ has compact support and thus vanishes at $\pm \infty$:
\[ \int_{-\infty}^{\infty} \frac{\partial a_i}{\partial x_i} dx_i = a_i(\dots, \infty, \dots) - a_i(\dots, -\infty, \dots) = 0 - 0 = 0. \]
}
Thus, $\int_M d\eta = 0$, which matches $\int_{\partial M} \eta = 0$.

\textbf{Stage 3: Case 2 - Support in a Boundary Chart}
Suppose the support of $\eta$ intersects the boundary $\partial M$. We use a coordinate chart $\phi: U \to \mathbb{H}^n$, where $\mathbb{H}^n = \{(x_1, \dots, x_n) \in \mathbb{R}^n : x_n \ge 0\}$. The boundary $\partial M \cap U$ corresponds to $x_n = 0$.
In these coordinates, the induced orientation on the boundary is $(-1)^n dx_1 \wedge \dots \wedge dx_{n-1}$.
\sidenote{
\textbf{Calculation for Boundary Support:}
1. \textbf{The Interior Integral:} As in Case 1, the integrals for $i < n$ vanish due to compact support. For $i=n$, we have $\eta_n = (-1)^{n-1} a_n dx_1 \wedge \dots \wedge dx_{n-1}$ and $d\eta_n = \frac{\partial a_n}{\partial x_n} dx_1 \wedge \dots \wedge dx_n$:
\[ \int_{\mathbb{H}^n} d\eta = \int_{\mathbb{R}^{n-1}} \left( \int_0^\infty \frac{\partial a_n}{\partial x_n} dx_n \right) dx_1 \dots dx_{n-1} \]
The inner integral is $a_n(x', \infty) - a_n(x', 0) = -a_n(x', 0)$. Thus:
\[ \int_M d\eta = -\int_{\mathbb{R}^{n-1}} a_n(x_1, \dots, x_{n-1}, 0) dx_1 \dots dx_{n-1}. \]
2. \textbf{The Boundary Integral:} On $x_n = 0$, $dx_n = 0$, so only the $i=n$ term survives:
\[ \eta|_{\partial \mathbb{H}^n} = (-1)^{n-1} a_n(x', 0) dx_1 \wedge \dots \wedge dx_{n-1}. \]
Using the boundary orientation form $\Omega_{\partial M} = (-1)^n dx_1 \wedge \dots \wedge dx_{n-1}$, we have $dx_1 \wedge \dots \wedge dx_{n-1} = (-1)^n \Omega_{\partial M}$.
\[ \int_{\partial M} \eta = \int_{\mathbb{R}^{n-1}} (-1)^{n-1} a_n(x', 0) (-1)^n \Omega_{\partial M} = -\int_{\mathbb{R}^{n-1}} a_n(x', 0) dx'. \]
Both integrals match, completing the proof.
}
\end{solution}

\clearpage

\begin{exercise}
\label{geometry:2012:12}
Let $M$ be a surface in $\mathbb{R}^{3}$. Let $D$ be a simply-connected domain in $M$ such that the boundary $\partial D$ is compact and consists of a finite number of smooth curves. Prove the Gauss-Bonnet Formula:
\[ \int_{\partial D} k_{g}\, ds + \sum_{j}(\pi - \alpha_{j}) + \iint_{D} K\, dA = 2\pi, \]
where $k_{g}$ is the geodesic curvature of the boundary curve. Each $\alpha_{j}$ is the interior angle at a vertex of the boundary, $K$ is the Gaussian curvature of $M$, and the 2-form $dA$ is the area element of $M$.
\end{exercise}

\begin{solution}
\textbf{Geometric Intuition:} The Gauss-Bonnet formula is a fundamental link between the local geometry of a surface and its global topology. For a simply-connected domain $D$, which is topologically equivalent to a disk, the formula relates the total Gaussian curvature within the domain to the "total turning" of its boundary. Imagine a bug traversing the boundary $\partial D$. The bug's direction changes continuously along smooth segments (measured by geodesic curvature) and abruptly at vertices (measured by exterior angles). On a flat plane, these changes must sum to $2\pi$ for the bug to return to its starting orientation. On a curved surface, the interior Gaussian curvature "compensates" for this turning; a positively curved surface (like a sphere) adds to the turning, while a negatively curved surface (like a saddle) subtracts from it.

\textbf{Proof:}
We proceed by using the method of moving frames. Since $D$ is a simply-connected domain on the surface $M$, we can choose a smooth orthonormal frame field $\{\vec{e}_1, \vec{e}_2\}$ over the entire domain $D$. Let $\{\theta^1, \theta^2\}$ be the corresponding dual coframe.

The Gaussian curvature $K$\footnote{The Gaussian curvature $K$ at a point on a surface is defined as the product of the principal curvatures, $K = \kappa_1 \kappa_2$. According to Gauss's \textit{Theorema Egregium}, $K$ is an intrinsic property of the surface, meaning it depends only on the first fundamental form and can be determined by measurements made within the surface.} is related to the connection 1-form $\omega_{12}$ by the second structural equation:
\[ d\omega_{12} = -K \theta^1 \wedge \theta^2 = -K dA, \]
where $dA$ is the area element of the surface. Applying Stokes' Theorem to the domain $D$, we obtain the relationship between the total Gaussian curvature and the line integral of the connection 1-form along the boundary:
\[ \iint_D K dA = -\iint_D d\omega_{12} = -\int_{\partial D} \omega_{12}. \]

Now, consider the boundary $\partial D$, which consists of piecewise smooth curves. Let $C(s)$ be a unit-speed parametrization of a smooth segment of $\partial D$. The tangent vector is given by $\dot{C}(s)$. Since $\{\vec{e}_1, \vec{e}_2\}$ is an orthonormal frame, we can express the tangent vector as:
\[ \dot{C}(s) = \cos\phi(s) \vec{e}_1 + \sin\phi(s) \vec{e}_2, \]
where $\phi(s)$ is the angle between the tangent vector and the frame vector $\vec{e}_1$. The geodesic curvature $k_g$\footnote{The geodesic curvature $k_g$ measures how much a curve deviates from being a geodesic within the surface. For a unit-speed curve $C(s)$, it is the projection of the acceleration vector $\ddot{C}$ onto the tangent plane of the surface, specifically in the direction perpendicular to the curve's tangent.} can be expressed in terms of the rate of change of this angle and the connection 1-form. Specifically, we have $k_g = \frac{d\phi}{ds} + \omega_{12}(\dot{C})$\sidenote{Let $\dot{C} = \cos\phi \vec{e}_1 + \sin\phi \vec{e}_2$. The covariant derivative is $\nabla_{\dot{C}} \dot{C} = (-\sin\phi \dot{\phi} \vec{e}_1 + \cos\phi \dot{\phi} \vec{e}_2) + (\cos\phi \nabla_{\dot{C}} \vec{e}_1 + \sin\phi \nabla_{\dot{C}} \vec{e}_2)$. Using $\nabla_{\dot{C}} \vec{e}_1 = \omega_{12}(\dot{C}) \vec{e}_2$ and $\nabla_{\dot{C}} \vec{e}_2 = -\omega_{12}(\dot{C}) \vec{e}_1$, we have $\nabla_{\dot{C}} \dot{C} = (-\sin\phi \dot{\phi} - \sin\phi \omega_{12}(\dot{C})) \vec{e}_1 + (\cos\phi \dot{\phi} + \cos\phi \omega_{12}(\dot{C})) \vec{e}_2$. The geodesic curvature is $k_g = \langle \nabla_{\dot{C}} \dot{C}, \vec{n} \times \dot{C} \rangle$. With $\vec{n} \times \dot{C} = -\sin\phi \vec{e}_1 + \cos\phi \vec{e}_2$, we find $k_g = \dot{\phi} + \omega_{12}(\dot{C})$.}.

Integrating $k_g$ over each smooth segment $C_i$ of the boundary gives:
\[ \int_{C_i} k_g ds = \int_{C_i} \frac{d\phi}{ds} ds + \int_{C_i} \omega_{12} = \Delta \phi_i + \int_{C_i} \omega_{12}, \]
where $\Delta \phi_i$ is the change in the angle $\phi$ along the segment $C_i$. Summing over all smooth segments that form the boundary $\partial D$:
\[ \int_{\partial D} k_g ds = \sum_i \Delta \phi_i + \int_{\partial D} \omega_{12} = \sum_i \Delta \phi_i - \iint_D K dA. \]

Rearranging the terms, we have:
\[ \int_{\partial D} k_g ds + \iint_D K dA = \sum_i \Delta \phi_i. \]

The boundary $\partial D$ is a piecewise smooth simple closed curve. At each vertex where two smooth segments $C_j$ and $C_{j+1}$ meet, there is a jump in the tangent direction. The exterior angle $\pi - \alpha_j$\footnote{The exterior angle $\epsilon_j = \pi - \alpha_j$ at a vertex is the angle through which the tangent vector rotates as it passes through the vertex. Here, $\alpha_j$ is the interior angle of the domain at that vertex.} represents this discrete rotation. According to the \textbf{Theorem of Turning Tangents} for a simply-connected domain with a global frame, the total continuous rotation plus the sum of discrete rotations at the vertices must equal $2\pi$:
\[ \sum_i \Delta \phi_i + \sum_j (\pi - \alpha_j) = 2\pi. \]
Substituting $\sum_i \Delta \phi_i$ from our previous result into this topological constraint, we obtain:
\[ \left( \int_{\partial D} k_g ds + \iint_D K dA \right) + \sum_j (\pi - \alpha_j) = 2\pi. \]
Rearranging the terms yields the Gauss-Bonnet formula for a simply-connected domain:
\[ \int_{\partial D} k_g ds + \sum_j (\pi - \alpha_j) + \iint_D K dA = 2\pi. \]
\end{solution}
